\chapter{2006年北京卷（理）}
\section{单选题}
\begin{problem}
在复平面内，复数 \(\frac{1+\i}{\i}\) 对应的点位于（\quad）

\choices
    {第一象限}
    {第二象限}
    {第三象限}
    {第四象限}
\end{problem}

\begin{answer}
D
\end{answer}

\begin{solution}
因为 \(\dfrac{1}{\i}=-\i\)，所以 \(\dfrac{1+\i}{\i}=1-\i\).
该复数的实部为正，虚部为负，故它在第四象限，答案为 D.
\end{solution}

\begin{problem}
若 \(\bs{a}\) 与 \(\bs{b}-\bs{c}\) 都是非零向量，则“\(\bs{a}\cdot\bs{b}=\bs{a}\cdot\bs{c}\)”是“\(\bs{a}\perp(\bs{b}-\bs{c})\)”的（\quad）

\choices
    {充分而不必要条件}
    {必要而不充分条件}
    {充分必要条件}
    {既不充分也不必要条件}
\end{problem}

\begin{answer}
C
\end{answer}

\begin{solution}
由数量积的分配律，\(\bs{a}\cdot(\bs{b}-\bs{c})=\bs{a}\cdot\bs{b}-\bs{a}\cdot\bs{c}\).
因此，\(\bs{a}\cdot\bs{b}=\bs{a}\cdot\bs{c}\) 当且仅当 \(\bs{a}\cdot(\bs{b}-\bs{c})=0\). 由于两向量均非零，这又当且仅当 \(\bs{a}\perp(\bs{b}-\bs{c})\)，所以两者互为充分必要条件，答案为 C.
\end{solution}

\begin{problem}
在 \(1\)，\(2\)，\(3\)，\(4\)，\(5\) 这五个数字组成的没有重复数字的三位数中，各位数字之和为奇数的共有（\quad）

\choices
    {\(36\)个}
    {\(24\)个}
    {\(18\)个}
    {\(6\)个}
\end{problem}

\begin{answer}
B
\end{answer}

\begin{solution}
五个数字中有奇数 \(1\)，\(3\)，\(5\) 共 \(3\) 个，偶数 \(2,4\) 共 \(2\) 个. 三位数的各位数字之和为奇数，当且仅当所取的奇数个数为奇数.

取三个奇数时，有 \(3!\) 个排法；取一个奇数和两个偶数时，有 \(3\times3!\) 个排法. 因此共有 \(3!+3\times3!=6+18=24\) 个，答案为 B.
\end{solution}

\begin{problem}
平面 \(\alpha\) 的斜线 \(AB\) 交 \(\alpha\) 于点 \(B\)，过定点 \(A\) 的动直线 \(l\) 与 \(AB\) 垂直，且交 \(\alpha\) 于点 \(C\)，则动点 \(C\) 的轨迹是（\quad）

\choices
    {一条直线}
    {一个圆}
    {一个椭圆}
    {双曲线的一支}
\end{problem}

\begin{answer}
A
\end{answer}

\begin{solution}
过点 \(A\) 且垂直于 \(AB\) 的所有直线都位于同一个平面 \(\beta\) 内，其中 \(\beta\) 过 \(A\) 且垂直于直线 \(AB\). 动点 \(C\) 同时满足 \(C\in\alpha\) 和 \(C\in\beta\)，故其轨迹是两平面 \(\alpha\)、\(\beta\) 的交线，即一条直线，答案为 A.
\end{solution}

\begin{problem}
已知 \(f(x)=\begin{cases}
(3a-1)x+4a, & x<1,\\
\log_{a}x, & x\geqslant1,
\end{cases}\) 是 \((-\infty,+\infty)\) 上的减函数，那么 \(a\) 的取值范围是（\quad）

\choices
    {\((0,1)\)}
    {\(\left(0,\frac{1}{3}\right)\)}
    {\(\left[\frac{1}{7},\frac{1}{3}\right)\)}
    {\(\left[\frac{1}{7},1\right)\)}
\end{problem}

\begin{answer}
C
\end{answer}

\begin{solution}
要使 \(f(x)\) 在整个实数范围内有定义，必须有 \(a>0\) 且 \(a\ne1\). 在 \(x\geqslant1\) 时，\(\log_a x\) 为减函数，故 \(0<a<1\). 在 \(x<1\) 时，线性函数的斜率必须为负，故
\(3a-1<0\)，
即 \(a<\frac13\). 此外，对任意 \(x<1\)，应有 \(f(x)>f(1)=0\). 左段在 \(x\to1^{-}\) 时的极限为 \(7a-1\)，所以必须有 \(7a-1\geqslant0\)，即 \(a\geqslant\frac17\). 综上，
\(\frac17\leqslant a<\frac13\)，
答案为 C.
\end{solution}

\begin{problem}
在下列四个函数中，满足性质“对于区间 \((1,2)\) 上的任意 \(x_1\)，\(x_2\)（\(x_1\neq x_2\)），\(|f(x_1)-f(x_2)|<|x_2-x_1|\) 恒成立”的只有（\quad）

\choices
    {\(f(x)=\frac{1}{x}\)}
    {\(f(x)=|x|\)}
    {\(f(x)=2^x\)}
    {\(f(x)=x^2\)}
\end{problem}

\begin{answer}
A
\end{answer}

\begin{solution}
对 \(f(x)=\dfrac1x\)，由中值定理，存在 \(\xi\in(1,2)\) 使
\[
\left|f(x_1)-f(x_2)\right|=\left|f'(\xi)\right|\left|x_1-x_2\right|=\frac{1}{\xi^2}\left|x_1-x_2\right|<\left|x_1-x_2\right|
\]
对 \(f(x)=|x|\)，因区间 \((1,2)\) 内 \(|x|=x\)，两边恒等，不能满足严格不等式. 对 \(f(x)=2^x\)，其导数在该区间内为 \(2^x\ln2>2\ln2>1\)；对 \(f(x)=x^2\)，其导数为 \(2x>2>1\)，后二者均不能满足条件. 因此只有 A.
\end{solution}

\begin{problem}
设 \(f(n)=2+2^{4}+2^{7}+2^{10}+\cdots+2^{3n+10}\)（\(n\in\N\)），则 \(f(n)\) 等于（\quad）

\choices
    {\(\frac{2}{7}(8^n-1)\)}
    {\(\frac{2}{7}(8^{n+1}-1)\)}
    {\(\frac{2}{7}(8^{n+3}-1)\)}
    {\(\frac{2}{7}(8^{n+4}-1)\)}
\end{problem}

\begin{answer}
D
\end{answer}

\begin{solution}
将各项写成等比数列，首项为 \(2\)，公比为 \(2^3=8\). 指数从 \(1\) 开始，以公差 \(3\) 递增到 \(3n+10=3(n+3)+1\)，所以共有 \(n+4\) 项. 因此 \(f(n)=2\frac{8^{n+4}-1}{8-1}=\frac27\left(8^{n+4}-1\right)\)，
答案为 D.
\end{solution}

\begin{problem}
如图为某三岔路口交通环岛的简化模型，在某高峰时段，单位时间进出路口 \(A\)，\(B\)，\(C\) 的机动车辆数如图所示，图中 \(x_1\)，\(x_2\)，\(x_3\) 分别表示该时段单位时间通过路段 \(\overline{AB}\)，\(\overline{BC}\)，\(\overline{CA}\) 的机动车辆数（假设：单位时间内，在上述路段中，同一路段上驶入与驶出的车辆数相等），则 \(x_1\)，\(x_2\)，\(x_3\) 的大小关系是（\quad）

\texfigure[max width=0.42\linewidth]{img_repaint/2006/beijing_science/q8_fig1.tex}

\choices
    {\(x_1>x_2>x_3\)}
    {\(x_1>x_3>x_2\)}
    {\(x_2>x_3>x_1\)}
    {\(x_3>x_2>x_1\)}
\end{problem}

\begin{answer}
C
\end{answer}

\begin{solution}
按图中箭头方向列各路口的流量守恒方程. 在路口 \(A\)，有 \(x_3+50=x_1+55\)，故 \(x_3=x_1+5\). 在路口 \(B\)，有 \(x_1+30=x_2+20\)，故 \(x_2=x_1+10\). 于是 \(x_2>x_3>x_1\)，
答案为 C.
\end{solution}

\section{填空题}
\begin{problem}
\(\displaystyle\lim_{x\to-1}\frac{x^2+3x+2}{x^2-1}\) 的值等于\fillinblank{}.
\end{problem}

\begin{answer}
\(-\dfrac12\)
\end{answer}

\begin{solution}
因式分解分子和分母，得
\[
\lim_{x\to-1}\frac{x^2+3x+2}{x^2-1}
=\lim_{x\to-1}\frac{(x+1)(x+2)}{(x+1)(x-1)}
=\lim_{x\to-1}\frac{x+2}{x-1}=-\frac12
\]
\end{solution}

\begin{problem}
在 \(\left(\sqrt{x}-\frac{2}{x}\right)^7\) 的展开式中，\(x^2\) 的系数为\fillinblank{}.（用数字作答）
\end{problem}

\begin{answer}
\(-14\)
\end{answer}

\begin{solution}
展开式的通项中，若选取 \(\left(-\dfrac2x\right)^j\) 这一项，则其中 \(x\) 的幂为 \(\left(\sqrt{x}\right)^{7-j}\left(\frac1x\right)^j=x^{(7-j)/2-j}=x^{(7-3j)/2}\). 令 \((7-3j)/2=2\)，得 \(j=1\). 因此 \(x^2\) 的系数为 \(\mathrm C_7^1(-2)=-14\).
\end{solution}

\begin{problem}
若三点 \(A(2,2)\)，\(B(a,0)\)，\(C(0,b)\)（\(ab\neq0\)）共线，则 \(\frac{1}{a}+\frac{1}{b}\) 的值等于\fillinblank{}.
\end{problem}

\begin{answer}
\(\dfrac12\)
\end{answer}

\begin{solution}
因为 \(A(2,2)\)、\(B(a,0)\)、\(C(0,b)\) 共线，且 \(ab\ne0\)，直线 \(BC\) 的截距式为 \(\frac{x}{a}+\frac{y}{b}=1\). 将点 \(A(2,2)\) 代入，得 \(\frac2a+\frac2b=1\)，所以 \(\frac1a+\frac1b=\frac12\).
\end{solution}

\begin{problem}
在 \(\triangle ABC\) 中，若 \(\sin A:\sin B:\sin C=5:7:8\)，则 \(\angle B\) 的大小是\fillinblank{}.
\end{problem}

\begin{answer}
\(\dfrac{\pi}{3}\)
\end{answer}

\begin{solution}
由正弦定理，三角形的边长之比等于对应角的正弦之比，因此可设 \(a=5k\)，\(b=7k\)，\(c=8k\)，其中 \(k>0\).
由余弦定理，
\[
\cos B=\frac{a^2+c^2-b^2}{2ac}
=\frac{25k^2+64k^2-49k^2}{2\cdot5k\cdot8k}=\frac12
\]
由于 \(B\) 是三角形内角，故 \(B=\dfrac\pi3\).
\end{solution}

\begin{problem}
已知点 \(P(x,y)\) 的坐标满足条件 \(\begin{cases}
x+y\leqslant4,\\
y\geqslant x,\\
x\geqslant1,
\end{cases}\)
点 \(O\) 为坐标原点，那么 \(|PO|\) 的最小值等于\fillinblank{}，最大值等于\fillinblank{}.
\end{problem}

\begin{answer}
\(\sqrt{2}\)，\(\sqrt{10}\)
\end{answer}

\begin{solution}
由 \(x\geqslant1\) 且 \(y\geqslant x\)，有 \(x\geqslant1\)、\(y\geqslant1\)，从而 \(|PO|^2=x^2+y^2\geqslant1+1=2\).
当 \((x,y)=(1,1)\) 时满足全部条件，故最小值为 \(\sqrt2\).

可行域是由三条边界直线围成的三角形，三个顶点分别为 \((1,1)\)、\((2,2)\)、\((1,3)\). 凸函数 \(x^2+y^2\) 在该三角形上的最大值取在顶点处，三个顶点到原点的距离分别为 \(\sqrt2\)、\(2\sqrt2\)、\(\sqrt{10}\)，所以最大值为 \(\sqrt{10}\).
\end{solution}

\begin{problem}
已知 \(A\)，\(B\)，\(C\) 三点在球心为 \(O\)，半径为 \(R\) 的球面上，\(AC\perp BC\)，且 \(AB=R\). 那么 \(A\)，\(B\) 两点的球面距离为\fillinblank{}，球心到平面 \(ABC\) 的距离为\fillinblank{}.
\end{problem}

\begin{answer}
\(\dfrac{\pi R}{3}\)，\(\dfrac{\sqrt{3}}{2}R\)
\end{answer}

\begin{solution}
设平面 \(ABC\) 与球面的截圆圆心为 \(O_1\)，半径为 \(r\). 因为 \(\angle ACB=90^\circ\)，由圆周角定理，弦 \(AB\) 是该截圆的直径，所以 \(r=\frac{AB}{2}=\frac R2\). 在直角三角形 \(OO_1A\) 中，\(OA=R\)、\(O_1A=r=R/2\)，故 \(OO_1=\sqrt{R^2-\left(\frac R2\right)^2}=\frac{\sqrt3}{2}R\).

设 \(A\)、\(B\) 对应的球心角为 \(\theta\)，则弦长公式给出 \(AB=2R\sin\frac\theta2\). 代入 \(AB=R\)，且 \(0<\theta\leqslant\pi\)，得 \(\sin\dfrac\theta2=\dfrac12\)，所以 \(\theta=\dfrac\pi3\). 因此两点的球面距离为 \(R\theta=\dfrac{\pi R}{3}\).
\end{solution}

\section{解答题}
\begin{problem}
已知函数 \(f(x)=\frac{1-\sqrt{2}\sin\left(2x-\frac{\pi}{4}\right)}{\cos x}\).

\begin{enumerate}
\item 求 \(f(x)\) 的定义域；
\item 设 \(\alpha\) 是第四象限的角，且 \(\tan\alpha=-\frac{4}{3}\)，求 \(f(\alpha)\) 的值.
\end{enumerate}
\end{problem}

\begin{solution}
\begin{enumerate}
\item 分式有意义的必要条件是分母不为零，即 \(\cos x\ne0\). 因此
\(x\ne k\pi+\frac\pi2\)（\(k\in\Z\)），
故定义域为 \(\left\{x\mid x\ne k\pi+\dfrac\pi2,\ k\in\Z\right\}\).
\item 由 \(\tan\alpha=-\dfrac43\)，且 \(\alpha\) 是第四象限的角，得
\(\sin\alpha=-\frac45\)，\(\cos\alpha=\frac35\).
利用和角公式，
\(\sqrt2\sin\left(2\alpha-\frac\pi4\right)=\sin2\alpha-\cos2\alpha\).
所以
\(f(\alpha)=\frac{1-\sin2\alpha+\cos2\alpha}{\cos\alpha}=\frac{1-2\sin\alpha\cos\alpha+\cos^2\alpha-\sin^2\alpha}{\cos\alpha}\). 再因 \(1-2\sin\alpha\cos\alpha+\cos^2\alpha-\sin^2\alpha=2\cos\alpha(\cos\alpha-\sin\alpha)\)，且 \(\cos\alpha\ne0\)，故 \(f(\alpha)=2(\cos\alpha-\sin\alpha)=\frac{14}{5}\).
\end{enumerate}
\end{solution}

\begin{problem}
已知函数 \(f(x)=ax^3+bx^2+cx\) 在点 \(x_0\) 处取得极大值 \(5\)，其导函数 \(y=f'(x)\) 的图象经过点 \((1,0)\)，\((2,0)\)，如图所示，求：

\begin{enumerate}
\item 求 \(x_0\) 的值；
\item 求 \(a\)，\(b\)，\(c\) 的值.
\end{enumerate}

\texfigure[max width=0.28\linewidth]{img_repaint/2006/beijing_science/q16_fig1.tex}
\end{problem}

\begin{solution}
\begin{enumerate}
\item 由图象可知，\(f'(x)>0\) 在 \((-\infty,1)\) 和 \((2,+\infty)\) 上成立，\(f'(x)<0\) 在 \((1,2)\) 上成立. 因此 \(f(x)\) 在 \(x=1\) 处由增变减，取得极大值，故 \(x_0=1\).
\item 由图象，\(f'(1)=f'(2)=0\). 又 \(f'(x)=3ax^2+2bx+c\)，
且二次项系数为 \(3a\)，所以
\(f'(x)=3a(x-1)(x-2)=3ax^2-9ax+6a\). 比较系数得 \(2b=-9a\)、\(c=6a\). 由 \(f(1)=5\)，有 \(a+b+c=a-\frac92a+6a=\frac52a=5\)，
解得 \(a=2\)，进而 \(b=-9\)，\(c=12\).
\end{enumerate}
\end{solution}

\begin{problem}
如图，在底面为平行四边形的四棱锥 \(P-ABCD\) 中，\(AB\perp AC\)，\(PA\perp\) 平面 \(ABCD\)，且 \(PA=AB\)，点 \(E\) 是 \(PD\) 的中点.

\begin{enumerate}
\item 求证：\(AC\perp PB\)；
\item 求证：\(PB\parallel\) 平面 \(AEC\)；
\item 求二面角 \(E-AC-B\) 的大小.
\end{enumerate}

\texfigure[max width=0.45\linewidth]{img_repaint/2006/beijing_science/q17_fig1.tex}
\end{problem}

\begin{solution}
\begin{enumerate}
\item 因为 \(PA\perp\) 平面 \(ABCD\)，且 \(AC\) 在平面 \(ABCD\) 内，所以 \(AC\perp PA\). 又已知 \(AC\perp AB\)，而 \(PA\)、\(AB\) 是平面 \(PAB\) 内相交的两条直线，故 \(AC\perp\) 平面 \(PAB\). 由于 \(PB\) 在平面 \(PAB\) 内，所以 \(AC\perp PB\).
\item 连接 \(BD\)，设 \(O=AC\cap BD\)，再连接 \(EO\). 由平行四边形的性质，\(O\) 是 \(BD\) 的中点；又 \(E\) 是 \(PD\) 的中点，所以在三角形 \(PDB\) 中，\(EO\parallel PB\). 因为 \(E\) 不在底面 \(ABCD\) 内，平面 \(AEC\) 与底面 \(ABCD\) 的交线为 \(AC\)，而 \(B\notin AC\)，故 \(PB\) 不在平面 \(AEC\) 内. 又 \(EO\) 在平面 \(AEC\) 内，因此 \(PB\parallel\) 平面 \(AEC\).
\item 过 \(O\) 作直线 \(FOG\parallel AB\)，其中 \(F\in AD\)，\(G\in BC\). 由三角形 \(ACD\) 及平行四边形的中点性质，\(F\) 是 \(AD\) 的中点，\(G\) 是 \(BC\) 的中点. 因为 \(OG\parallel AB\perp AC\)，所以 \(OG\perp AC\)；由前两问的结论，\(EO\parallel PB\perp AC\)，所以 \(EO\perp AC\). 因而 \(\angle EOG\) 是二面角 \(E-AC-B\) 的平面角.

在三角形 \(PAD\) 中，\(E\)、\(F\) 分别是 \(PD\)、\(AD\) 的中点，所以 \(EF\parallel PA\)，且 \(EF=\frac12PA=\frac12AB\). 又由三角形 \(ACD\) 的中位线性质，\(OF\parallel CD\parallel AB\)，且 \(OF=\frac12AB\). 因为 \(PA\perp\) 平面 \(ABCD\)，所以 \(EF\perp OF\). 因此直角三角形 \(EFO\) 的两直角边相等，\(\angle EOF=45^\circ\). 由于 \(F\)、\(O\)、\(G\) 依次在同一直线上，\(\angle EOG=180^\circ-\angle EOF=135^\circ\). 所以二面角 \(E-AC-B\) 的大小为 \(135^\circ\).
\end{enumerate}
\end{solution}

\begin{problem}
某公司招聘员工，指定三门考试课程，有两种考试方案.
方案一：考试三门课程，至少有两门及格为考试通过；
方案二：在三门课程中，随机选取两门，这两门都及格为考试通过.
假设某应聘者对三门指定课程考试及格的概率分别是 \(a\)，\(b\)，\(c\)，且三门课程考试是否及格相互之间没有影响.

\begin{enumerate}
\item 分别求该应聘者用方案一和方案二时考试通过的概率；
\item 试比较该应聘者在上述两种方案下考试通过的概率的大小.（说明理由）
\end{enumerate}
\end{problem}

\begin{solution}
\begin{enumerate}
\item 记三门课程及格的事件分别为 \(A\)、\(B\)、\(C\)，则
\(P(A)=a\)，\(P(B)=b\)，\(P(C)=c\).
方案一通过包括恰有两门及格和三门都及格两种互斥情形，因此
\(p_1=ab(1-c)+ac(1-b)+bc(1-a)+abc=ab+ac+bc-2abc\). 方案二随机选取三组课程中的一组，每组被选中的概率为 \(\dfrac13\)，所以 \(p_2=\frac13\left(ab+ac+bc\right)\).
\item 因为 \(0\leqslant a\leqslant1\)，\(0\leqslant b\leqslant1\)，\(0\leqslant c\leqslant1\)，有
\(p_1-p_2=\frac23\left(ab+ac+bc-3abc\right)=\frac23\left[ab(1-c)+ac(1-b)+bc(1-a)\right]\geqslant0\).
故方案一的通过概率不小于方案二；当且仅当括号内的三个非负项均为零时取等号，否则方案一的概率严格较大.
\end{enumerate}
\end{solution}

\begin{problem}
已知点 \(M(-2,0)\)，\(N(2,0)\)，动点 \(P\) 满足条件 \(\left|PM\right|-\left|PN\right|=2\sqrt{2}\). 记动点 \(P\) 的轨迹为 \(W\).

\begin{enumerate}
\item 求 \(W\) 的方程；
\item 若 \(A\)，\(B\) 是 \(W\) 上的不同两点，\(O\) 是坐标原点，求 \(\overrightarrow{OA}\cdot\overrightarrow{OB}\) 的最小值.
\end{enumerate}
\end{problem}

\begin{solution}
\begin{enumerate}
\item 由双曲线的定义，\(W\) 是以 \(M\)、\(N\) 为焦点的双曲线的右支. 其焦距的一半为 \(c=2\)，实半轴长满足 \(2a=2\sqrt2\)，故 \(a=\sqrt2\). 于是 \(b^2=c^2-a^2=4-2=2\)，所以 \(W\) 的方程为 \(\frac{x^2}{2}-\frac{y^2}{2}=1\)，且 \(x\geqslant\sqrt2\).
\item 设 \(A(x_1,y_1)\)、\(B(x_2,y_2)\)，并令 \(s_i=x_i+y_i\)，\(t_i=x_i-y_i\)（\(i=1\)，\(2\)）.
由点 \(A\)、\(B\) 在双曲线右支上，得 \(s_i>0\)、\(t_i>0\)，且
\(s_it_i=x_i^2-y_i^2=2\).
因此
\(\overrightarrow{OA}\cdot\overrightarrow{OB}=x_1x_2+y_1y_2=\frac12\left(s_1s_2+t_1t_2\right)\).
由算术几何平均不等式，
\[
\overrightarrow{OA}\cdot\overrightarrow{OB}
\geqslant\sqrt{s_1s_2t_1t_2}=2
\]
其中不等式为算术几何平均不等式. 取 \(A\left(\dfrac32,\dfrac12\right)\)、\(B\left(\dfrac32,-\dfrac12\right)\)，两点均在 \(W\) 上且互不相同，并且内积等于 \(2\). 故所求最小值为 \(2\).
\end{enumerate}
\end{solution}

\begin{problem}
在数列 \(\{a_n\}\) 中，若 \(a_1\)，\(a_2\) 是正整数，且 \(a_n=\left|a_{n-1}-a_{n-2}\right|\)，\(n=3\)，\(4\)，\(5\)，\(\cdots\)，则称 \(\{a_n\}\) 为“绝对差数列”.

\begin{enumerate}
\item 举出一个前五项不为零的“绝对差数列”（只要求写出前十项）；
\item 若“绝对差数列” \(\{a_n\}\) 中，\(a_{20}=3\)，\(a_{21}=0\)，数列 \(\{b_n\}\) 满足 \(b_n=a_n+a_{n+1}+a_{n+2}\)，\(n=1\)，\(2\)，\(3\)，\(\cdots\)，分别判断当 \(n\to\infty\) 时，\(a_n\) 与 \(b_n\) 的极限是否存在，如果存在，求出其极限值；
\item 证明：任何“绝对差数列”中总含有无穷多个为零的项.
\end{enumerate}
\end{problem}

\begin{solution}
\begin{enumerate}
\item 取 \(a_1=1\)、\(a_2=3\)，依次计算得
\(a_3=2\)，\(a_4=1\)，\(a_5=1\)，\(a_6=0\)，\(a_7=1\)，\(a_8=1\)，\(a_9=0\)，\(a_{10}=1\).
因此，前十项可以写为 \(1\)，\(3\)，\(2\)，\(1\)，\(1\)，\(0\)，\(1\)，\(1\)，\(0\)，\(1\)，其前五项均不为零.
\item 由递推关系，
\(a_{22}=|a_{21}-a_{20}|=3\)，
进而
\(a_{23}=|a_{22}-a_{21}|=3\)，\(a_{24}=|a_{23}-a_{22}|=0\).
所以从第 \(20\) 项开始，数列按 \(3\)，\(0\)，\(3\) 为周期循环，即
\(a_{20+3k}=3\)，\(a_{21+3k}=0\)，\(a_{22+3k}=3\)（\(k=0\)，\(1\)，\(2\)，\(\ldots\)）.
因此 \(a_n\) 有值为 \(0\) 和 \(3\) 的两个子列，极限不存在. 另一方面，每三个连续项的和均为 \(6\)，故当 \(n\geqslant20\) 时 \(b_n=6\)，从而
\(\displaystyle\lim_{n\to\infty}b_n=6\).
\item 先证明数列中必定出现零项. 反设所有项都不为零，则由于各项为非负整数，所有项都为正整数. 令
\(c_k=\max\{a_{2k-1},a_{2k}\}\)（\(k=1\)，\(2\)，\(3\)，\(\ldots\)）.
在没有零项的假设下，\(a_{2k+1}\ne0\)，故 \(a_{2k}\ne a_{2k-1}\)，所以
\(a_{2k+1}=|a_{2k}-a_{2k-1}|<c_k\).
又因 \(a_{2k}>0\) 且 \(a_{2k+1}<c_k\)，有
\(a_{2k+2}=|a_{2k+1}-a_{2k}|<c_k\).
于是 \(0<c_{k+1}<c_k\)，这给出一个无限下降的正整数列，不可能成立. 因此绝对差数列必有某项为零.

设第一个零项为 \(a_m=0\)，则 \(a_{m-1}=A>0\). 由递推关系，
\(a_{m+1}=|a_m-a_{m-1}|=A\)，
\(a_{m+2}=|a_{m+1}-a_m|=A\)，
\(a_{m+3}=|a_{m+2}-a_{m+1}|=0\).
同样的计算可反复进行，得到
\(a_{m+3k}=0\)，\(a_{m+3k+1}=A\)，\(a_{m+3k+2}=A\)
对所有 \(k=0\)，\(1\)，\(2\)，\(\ldots\) 成立. 因此数列中含有无穷多个为零的项.
\end{enumerate}
\end{solution}
